
\section{菲涅耳半波带与半波带法}

\begin{definition}[菲涅耳半波带]
\textbf{惠更斯原理}说，狭缝 \(AB\) 处波阵面上的每一点都可以看成新的波源，向各个方向发射子波。对观察点 \(P\) 而言，把波阵面上到 \(P\) 的光程差按 \(\lambda/2\) 逐段划分，得到的一段一段区域就叫\textbf{菲涅耳半波带}。相邻半波带到 \(P\) 的相位差为 \(\pi\)，因此最容易成对抵消。

对单缝衍射而言，衍射角 \(\theta\) 是衍射光与狭缝法线之间的夹角；屏上某点的光强，就是狭缝波阵面上所有具有相同衍射角的子波在该点的相干叠加。不同的 \(\theta\) 或不同的观察点 \(P\)，对应的光程差不同，叠加强度也不同。狭缝上下边缘 \(A,B\) 之间的最大光程差为
\[
\delta=BC=a\sin\theta.
\]
当 \(\theta=0\) 时，\(\delta=0\)，中心 \(O\) 为明纹（主极大）。
\end{definition}

\begin{figure}[htbp]
\centering
\includegraphics[width=0.6\textwidth]{pdf-exercises/9.png}
\caption{半波带法示意图}
\end{figure}

\begin{definition}[半波带法]
缝上下边缘到 \(Q\) 的最大光程差：$\delta=a\sin\theta$
\begin{itemize}
\item 中央明纹中心
\[
a\sin\theta=0
\]所有子波到屏上中心点的光程差都为 \(0\)，同相叠加，所以是中央明纹中心，也是主极大。
\item 暗纹
\[
a\sin\theta=\pm 2k\frac{\lambda}{2}=\pm k\lambda
\]这时单缝可分成 \(2k\) 个半波带。相邻半波带两两抵消，光强为零，形成暗纹。
\item 次级明纹
\[
a\sin\theta=\pm(2k+1)\frac{\lambda}{2}
\]这时单缝可分成 \(2k+1\) 个半波带。前 \(2k\) 个成对抵消，剩下一个半波带，所以形成明纹。但它不是主极大，亮度比中央明纹弱很多。
\item 介于明暗之间
\[
a\sin\theta\neq k\frac{\lambda}{2}
\]这时不能刚好分成整数个半波带，抵消不完全，也不是严格的明纹中心，所以光强介于明暗之间。
\end{itemize}
\[
\boxed{
\begin{array}{c|c}
a\sin\theta=0 & \text{中央明纹中心}\\
a\sin\theta=\pm k\lambda & \text{暗纹}\\
a\sin\theta=\pm(2k+1)\lambda/2 & \text{次级明纹}\\
\text{其他} & \text{明暗之间}
\end{array}}
\]这里 \(k=1,2,3,\dots\)。注意中央明纹不能简单归入“奇数个半波带”，因为它对应的是光程差为零，不是分出一个半波带。
\end{definition}
